\section*{Capítulo \ref{ch:b3:complex-mechanisms} --- Mecanismos de reação dos complexos}

\begin{solution}{exo:b3:complex-mechanisms:1}
Inertes: $\ce{[Cr(H2O)6]^{3+}}$ ($d^3$) e $\ce{[Co(NH3)6]^{3+}}$ ($d^6$ de spin baixo), grandes energias de ativação
do campo ligante. \omterm{def:b3:complex-mechanisms:labile}{Lábeis}: $\ce{[Cr(H2O)6]^{2+}}$ ($d^4$ de spin alto) e $\ce{[Cu(H2O)6]^{2+}}$ ($d^9$),
distorcidos por Jahn--Teller, com ligações axiais longas e fracas; $\ce{[Zn(H2O)6]^{2+}}$ ($d^{10}$), sem barreira de
campo ligante.
\end{solution}

\begin{solution}{exo:b3:complex-mechanisms:2}
Quando $k_2[\mathrm Y] \gg k_{-1}[\mathrm X]$, o denominador é $k_2[\mathrm Y]$ e $v = k_1[\mathrm{ML_5X}]$; quando $k_{-1}[\mathrm X] \gg k_2[\mathrm Y]$, ele é
$k_{-1}[\mathrm X]$ e segue-se a segunda forma (um pré-equilíbrio inibido pelo ligante
de saída).
\end{solution}

\begin{solution}{exo:b3:complex-mechanisms:3}
D: $\Delta V^{\ddagger} > 0$ e $\Delta^{\ddagger}S^\circ > 0$ (um ligante sai). A: ambos negativos (um ligante se liga no
estado de transição).
\end{solution}

\begin{solution}{exo:b3:complex-mechanisms:4}
$cis$-$\ce{[PtCl2(NH3)2]}$ a partir de $\ce{[PtCl4]^{2-}}$; $trans$-$\ce{[PtCl2(NH3)2]}$ a partir de $\ce{[Pt(NH3)4]^{2+}}$ (o
\omterm{def:b3:complex-mechanisms:trans-effect}{efeito trans} de \ce{Cl-} supera o de \ce{NH3}).
\end{solution}

\begin{solution}{exo:b3:complex-mechanisms:5}
$k_{\mathrm{obs}} = 2.0 \times \num{3.0e4} \times 0.10/1.2 = \qty{5.0e3}{s^{-1}}$; a \qty{1.0}{mol/L}, $\num{6.0e4}/3.0 = \qty{2.0e4}{s^{-1}}$; limite
$k_i = \qty{3.0e4}{s^{-1}}$.
\end{solution}

\begin{solution}{exo:b3:complex-mechanisms:6}
$\Delta V^{\ddagger} = -RT\ln2/\Delta p = -8.314 \times 298 \times 0.693/\num{99.9e6} = \qty{-1.7e-5}{m^3.mol^{-1}} = \qty{-17}{cm^3.mol^{-1}}$:
associativo.
\end{solution}

\begin{solution}{exo:b3:complex-mechanisms:7}
A fosfina, forte $\sigma$-doadora e $\pi$-aceptora, enfraquece a ligação em trans a ela (\omterm{def:b3:complex-mechanisms:trans-effect}{influência
trans}, ligação mais longa) e a labiliza (\omterm{def:b3:complex-mechanisms:trans-effect}{efeito trans}): esse cloreto é
substituído primeiro.
\end{solution}

\begin{solution}{exo:b3:complex-mechanisms:8}
O ligante em ponte acaba transferido para o metal oxidado; a velocidade depende
fortemente da natureza da ponte potencial; detecta-se um intermediário binuclear;
e a velocidade não pode exceder a da substituição necessária para formar a
ponte no parceiro \omterm{def:b3:complex-mechanisms:labile}{lábil}.
\end{solution}

\begin{solution}{exo:b3:complex-mechanisms:9}
$(\lambda + \Delta G^\circ)^2/4\lambda$ com $4\lambda = 320$: 20.0, 11.3, 0 e \qty{5.0}{kJ/mol} (o último na \omterm{def:b3:complex-mechanisms:inverted}{região
invertida}).
\end{solution}

\begin{solution}{exo:b3:complex-mechanisms:10}
$k_{12} = \sqrt{4.0 \times \num{2.0e5} \times \num{1.0e4}} = \qty{8.9e4}{L.mol^{-1}.s^{-1}}$.
\end{solution}

\begin{solution}{exo:b3:complex-mechanisms:11}
$d^5$ de spin alto: 0; $d^7$: $1.33\,Dq$; $d^8$: $2.0\,Dq$. Velocidades de troca de água esperadas: \ce{Mn^{2+}} $>$ \ce{Co^{2+}} $>$
\ce{Ni^{2+}}, a ordem observada.
\end{solution}

\begin{solution}{exo:b3:complex-mechanisms:12}
Para reações muito exergônicas, a velocidade intrínseca excede a velocidade com que os íons
se encontram: a constante observada fica no limite de difusão e a diminuição fica
escondida. Com doador e aceptor fixos a uma distância definida em uma só molécula, não
há etapa de difusão, e uma série de aceptores de força motriz crescente
mostrou a velocidade subindo, passando por um máximo e caindo.
\end{solution}

\begin{solution}{pb:b3:complex-mechanisms:1}
\textbf{1.} $5d^8$.
\textbf{2.} O campo forte de um íon $5d$ deixa o orbital $d_{x^2-y^2}$, que aponta para os quatro
ligantes, muito alto: os oito elétrons se emparelham nos quatro orbitais mais baixos.
\textbf{3.} O complexo tem 16 elétrons de valência e um sítio axial livre: o ligante de
entrada se liga primeiro, por um estado de transição pentacoordenado.
\textbf{4.} $k_{\mathrm{obs}} = k_1 + k_2[\mathrm Y]$: um caminho do solvente (ataque do solvente, depois substituição rápida por Y)
e um caminho direto.
\textbf{5.} A substituição leva horas (Parte III): inerte na escala dos minutos, não na de
um dia.
\textbf{6.} Cisplatina, porque o segundo \ce{NH3} substitui um cloreto em trans a um cloreto
(\omterm{def:b3:complex-mechanisms:trans-effect}{efeito trans} $\ce{Cl-} > \ce{NH3}$); também se formam subprodutos.
\textbf{7.} O iodeto tem um \omterm{def:b3:complex-mechanisms:trans-effect}{efeito trans} maior que o cloreto, o que torna o efeito
diretor mais limpo e mais rápido.
\textbf{8.} No $\ce{[PtI3(NH3)]-}$, é o iodeto em trans a um iodeto que é labilizado, não o que está em trans ao
\ce{NH3}: a segunda amônia entra em cis à primeira.
\textbf{9.} O iodeto de prata precipita e o $cis$-$\ce{[Pt(H2O)2(NH3)2]^{2+}}$ permanece em solução.
\textbf{10.} O cloreto substitui as duas moléculas de água em suas próprias posições; os
ligantes amin não são tocados.
\textbf{11.} A partir de $\ce{[Pt(NH3)4]^{2+}}$ e dois cloretos.
\textbf{12.} Suas duas posições reativas são opostas: ela não pode se ligar a duas
bases vizinhas da mesma fita de DNA, a lesão pela qual a cisplatina
age.
\textbf{13.} $\ln2/\num{9.0e-5} = \qty{7.7e3}{s} = \qty{2.1}{h}$.
\textbf{14.} $1 - \eu^{-\num{9.0e-5} \times 7200} = 0.48$.
\textbf{15.} A água é um grupo de saída muito melhor que o cloreto: o aquacomplexo é
substituído rapidamente pelo nitrogênio de uma base do DNA.
\textbf{16.} Associativa, como todas as substituições da platina(II): \omterm{def:b3:complex-mechanisms:activation-volume}{volume de ativação}
negativo.
\textbf{17.} A alta concentração de cloreto desloca o equilíbrio de volta para a
forma dicloro, mantendo o fármaco intacto até que ele chegue às células.
\textbf{18.} $f = K/(K + [\ce{Cl-}])$.
\textbf{19.} $\num{3.0e-3}/0.103 = 0.029$.
\textbf{20.} $\num{3.0e-3}/0.0070 = 0.43$.
\textbf{21.} 15.
\textbf{22.} Só dentro das células uma fração apreciável existe na forma de aquacomplexo
reativo.
\textbf{23.} Os ligantes de enxofre, como a glutationa e a cisteína e a metionina das
proteínas, que se ligam fortemente à platina, um metal mole.
\textbf{24.} A amônia se liga fortemente à platina e fica em trans ao cloreto, um ligante de
pequeno \omterm{def:b3:complex-mechanisms:trans-effect}{efeito trans}: as ligações com os ligantes amin não são labilizadas.
\textbf{25.} \textbf{Cerca de 15 vezes mais fármaco está na forma aqua reativa dentro de uma
célula do que no plasma.}
\end{solution}
